Este trabajo nació de una pregunta concreta: ¿es posible que la inteligencia artificial ayude a aprender a programar sin quitarle al alumno el esfuerzo de hacerlo?

La pregunta tiene más miga de lo que parece. No se trata de si la IA puede generar código, porque eso ya lo hace, y lo hace bien. Se trata de si puede hacer lo contrario, estar presente en el proceso de aprendizaje sin cortocircuitarlo. Negarse a resolver. Preguntar en lugar de responder. Acompañar sin sustituir.

SocratiCode es el intento de responder esa pregunta con software. Lo que empezó como una idea sobre \textit{prompt engineering} terminó siendo una plataforma completa con editor de código, ejecución en entorno aislado, moderación en tiempo real y una interfaz pensada para que el alumno y el tutor compartan el mismo espacio de trabajo. Cada decisión técnica tuvo detrás una decisión pedagógica.

Espero que este documento, y el sistema que describe, sean útiles para quien quiera construir algo parecido, o simplemente para quien le haya dado vueltas a la misma pregunta.

\hfill \begin{flushright}Ignacio García Hernanz\end{flushright}\hfill
